\section{Limitations}
\label{sec:09}
\label{sec:limitations}

VoteChain-ZK relies on identity-aware biometric enrollment and credential
authorization. Although the authorization request omits candidate choice, the
credential custodian can recover each stored voting secret, derive its
election-scoped nullifier, and associate a public ballot transaction with an
identified voter. This identity link does not itself reveal the encrypted candidate
selection. Moreover, the Groth16 membership proof does not establish one credential
per eligible person or verify that a fresh biometric match occurred when the
transaction was submitted.

Ballot confidentiality also depends on the voting client and trustee-share custody.
A compromised browser can observe the plaintext selection before encryption; the
circuit cannot attest to client integrity or the quality of locally generated
randomness. Three colluding trustees with sufficient shares could decrypt individual
published ciphertexts as well as aggregate totals, and independent organizational
custody of the five shares was not established. The implementation does not
establish receipt-freeness or coercion resistance.

The evaluated Groth16 setup used development-stage artifacts rather than an
independently audited production ceremony. Ballot proofs are checked on-chain,
whereas distributed-key-generation records and trustee discrete-logarithm equality
proofs are verified externally. The contract's finalization checks enforce election
state, padded-slot restrictions, count bounds, and agreement with the
accepted-ballot count; they do not prove that published totals were correctly
decrypted. Independent auditing requires the complete finalization transcript and
verification against the on-chain commitments and encrypted aggregates.

Experiments used fixed Merkle depth 17 and 16 encrypted ballot slots.
Candidate-count trials varied active candidates, not circuit width, while
eligible-set trials used synthetic commitments and fixed-depth membership paths.
Consequently, the results do not establish costs for deeper circuits, greater
ballot-slot capacity, or national-scale enrollment. Measured proof-stage intervals
include runtime overhead and should not be treated as isolated
cryptographic-operation timings.

Transaction measurements came from a local Anvil network and cover
submission-to-receipt intervals, not public-network confirmation, throughput, or
actual fees. They exclude biometric acquisition, authorization, credential delivery,
and proof preparation; the predecessor system's biometric evaluation was not
repeated. Complete threshold finalization also lacks a repeated timing campaign.
Production readiness requires stronger client and credential protections,
independent trustee administration, adversarial and concurrent-load testing,
reliable transcript availability, operational recovery procedures, and end-to-end
voter-session measurements.
